\unitopen{43}
\begin{candos}
\item przeczytać szwajcarską umowę o pracę i znaleźć w niej kluczowe punkty (wynagrodzenie, okres próbny, wypowiedzenie, urlop);
\item negocjować wynagrodzenie i warunki po niemiecku, uprzejmie i konkretnie;
\item odpowiedzieć na ofertę (przyjąć, odrzucić, złożyć kontrofertę);
\item znać podstawowe pojęcia płacowe w Szwajcarii (brutto, netto, 13.\ pensja, \de{Pensionskasse}).
\end{candos}

\begin{uwaga}[Nie jest to porada prawna]
Wszystkie twierdzenia o prawie pracy w tej jednostce oznaczone \verify{w ten sposób} wymagają sprawdzenia w tekście ustawy
(Obligationenrecht), w układzie zbiorowym i w Twojej umowie. Nie podejmuj decyzji na podstawie tej książki.
\end{uwaga}

\section{Słownictwo}
\vocabtable{43}

\section{Elementy umowy o pracę w Szwajcarii}
\begin{center}\small
\begin{tabular}{@{}p{0.24\linewidth}p{0.68\linewidth}@{}}
\toprule
\textsf{\textbf{Element}} & \textsf{\textbf{Co sprawdzasz}}\\
\midrule
\de{Stellenantritt} & data rozpoczęcia pracy\\
\de{Pensum} & wymiar etatu (\de{100\%}, \de{80\%})\\
\de{Lohn} & wynagrodzenie roczne brutto, \de{13.\ Monatslohn} \verify{(zwyczajowe, ale nie ustawowe)}, \de{Bonus}\\
\de{Probezeit} & okres próbny \verify{(ustawowo domyślnie jeden miesiąc, maksymalnie trzy, art.\ 335b OR)}\\
\de{Kündigungsfrist} & wypowiedzenie po okresie próbnym \verify{(domyślnie miesiąc w pierwszym roku, dwa w latach 2--9, trzy od dziesiątego roku, art.\ 335c OR)}\\
\de{Ferien} & urlop \verify{(minimum cztery tygodnie rocznie, art.\ 329a OR)}\\
\de{Überstunden} & zasady rekompensaty \verify{(zależy od umowy; art.\ 321c OR)}\\
\de{Pensionskasse} & ubezpieczenie emerytalne pracownicze \verify{(BVG)}\\
\de{Geheimhaltung} & klauzula poufności; czasem \de{Konkurrenzverbot}\\
\bottomrule
\end{tabular}
\end{center}

\section{Negocjacje}
Zwroty uprzejme i konkretne: \de{Ich würde mich über ein Gehalt von X freuen.} \de{Das entspricht nicht ganz meinen Vorstellungen.}
\de{Könnten wir uns bei Y treffen?} \de{Wären Sie bereit, den Bonus zu erhöhen?} \de{Ich bin offen für ein Gespräch zum Gesamtpaket.}
Odpowiedzi na ofertę: \de{Vielen Dank für das Angebot, ich nehme es gern an.} \de{Ich muss das Angebot noch prüfen und melde mich bis Freitag.}
\de{Leider muss ich das Angebot ablehnen.}

\begin{zasada}[Negocjacje po szwajcarsku]
Szwajcarski styl to rzeczowość, fakty i spokój. Zbierz dane (zakresy płac, koszty życia), mów o \emph{wartości}, którą dajesz, i pytaj o
\emph{pakiet} (urlop, edukacja, godziny, praca zdalna), nie tylko o pensję. Nie blefuj ani nie wywieraj presji: unika się tego.
\end{zasada}

\begin{falle}
Polski kontekst skłania do skromności albo odwrotnie do nacisku; w Szwajcarii odbiera się to źle. Błąd językowy: \de{Ich will mehr Geld} (zbyt
bezpośrednie). Poprawnie: \de{Ich würde mir ein Gehalt von X vorstellen.} Drugi: pomylenie \de{brutto} i \de{netto} w rozmowie; zawsze
pytaj \de{brutto oder netto?} i \de{pro Jahr oder pro Monat?}. Trzeci: nie podpisuj umowy, której nie rozumiesz; poproś o czas do namysłu.
\end{falle}

\begin{szwajcaria}
Wynagrodzenie podaje się zwykle w CHF brutto rocznie, a składki (\de{AHV}, \de{ALV}, \de{BVG}, \de{UVG}) pomniejszają kwotę netto. Pracownik zagraniczny
bywa objęty \de{Quellensteuer} \verify{(zależy od pozwolenia, kantonu i dochodu)}. Poproś o wyliczenie netto albo użyj kalkulatora kantonalnego.
\end{szwajcaria}

\begin{praca}
Odpowiedź na ofertę: \de{Sehr geehrte Frau Keller, vielen Dank für Ihr Angebot. Ich freue mich über die Stelle. Bevor ich zusage, hätte ich
gern ein Gespräch über das Gehalt und die Kündigungsfrist.} Wyślij takie pytanie e-mailem, nie zakładaj nic ustnie.
\end{praca}

\section{Dialogi}
\textbf{Dialog 1 --- negocjacje}
\begin{dialog}
\turn{A}{Wir bieten Ihnen ein Jahresgehalt von 125.000 Franken brutto an.}{Oferujemy roczne wynagrodzenie 125 000 franków brutto.}
\turn{B}{Vielen Dank. Das liegt etwas unter meinen Vorstellungen. Ich hatte mit 140.000 gerechnet.}{Dziękuję. To trochę poniżej moich oczekiwań. Liczyłem na 140 000.}
\turn{A}{Wir könnten 130.000 anbieten, dazu einen Bonus und fünf Wochen Ferien.}{Moglibyśmy zaproponować 130 000, do tego premię i pięć tygodni urlopu.}
\turn{B}{Das klingt gut. Könnten Sie mir das schriftlich schicken, damit ich es prüfen kann?}{Brzmi dobrze. Czy mogłaby mi Pani to przesłać na piśmie, żebym mógł to sprawdzić?}
\end{dialog}

\section{Ćwiczenia}
\begin{cwiczenia}
\exitem{Zapytaj, czy wynagrodzenie jest brutto czy netto.}{\de{Ist das Gehalt brutto oder netto?}}
\exitem{Zapytaj o okres wypowiedzenia.}{\de{Wie lang ist die Kündigungsfrist?}}
\exitem{Zapytaj o liczbę dni urlopu.}{\de{Wie viele Ferientage habe ich pro Jahr?}}
\exitem{Uprzejmie zasygnalizuj niezadowolenie z oferty.}{\de{Das entspricht nicht ganz meinen Vorstellungen.}}
\exitem{Zaproponuj kontrofertę (kwota X).}{\de{Ich würde mir ein Gehalt von X Franken vorstellen.}}
\exitem{Poproś o ofertę na piśmie.}{\de{Könnten Sie mir das Angebot bitte schriftlich schicken?}}
\exitem{Poproś o czas do namysłu.}{\de{Ich möchte das Angebot in Ruhe prüfen und melde mich bis Freitag.}}
\exitem{Przyjmij ofertę.}{\de{Vielen Dank, ich nehme das Angebot gern an.}}
\exitem{Odrzuć ofertę uprzejmie.}{\de{Leider muss ich das Angebot ablehnen, danke für Ihr Interesse.}}
\exitem{Przetłumacz: „Umowa jest na czas nieokreślony z okresem próbnym trzech miesięcy.”}{\de{Der Vertrag ist unbefristet mit einer Probezeit von drei Monaten.}}
\end{cwiczenia}

\section{Trening słownictwa (generowany)}
\drill{43}

\weekend{Przeczytaj wzór szwajcarskiej umowy o pracę (strona kantonalna lub urzędowa) i wypisz po niemiecku dziesięć pojęć z tłumaczeniem.
Zagraj z AI scenę negocjacji (ty, rekruter) i nagraj ją. Sprawdź tylko uprzejmość i konkret, nie liczby.}
